% included by aGMX_clustering.tex: phase 2 of 26 October 2026, the exact pixel integrals of the cloud Gaussians.
% Numbers from the validation runs named in each paragraph (clustering_tests/gmx.py, PIX_EXACT).
%======================================================================================================
\section{Exact pixel integrals (phase 2, 26 October)}\label{sec:phase2}

\paragraph{The approximation it replaces.} Step~2 evaluates each cloud at the pixel centre with its covariance increased by
$\tfrac1{12}$ (the variance of a uniform distribution on the pixel), and the M-step subtracts $\tfrac1{12}$ again. This is
exact for the second moment but not for the profile: for a cloud with $\sigma\lesssim0.3$\,px almost all the charge is in
one or two pixels, the approximated peak is too high and the tails too low, and the EM compensates by moving charge
between neighbouring clouds. The visible symptom was a sharp (front-side) X-ray next to a diffuse one: the sharp one was
lost or merged in 22\,\% of cases and its energy was within 5\,\% only 35\,\% of the time.

\paragraph{The exact model.} A cloud with covariance $\Sigma_k$ deposits in pixel $p=(i,j)$ (pixel centres at integers)
the charge
\begin{equation}
g_k(p)=A_k\int_{i-\frac12}^{i+\frac12}\!\!\int_{j-\frac12}^{j+\frac12}\N\!\big(\mathbf x;\bm\mu_k,\Sigma_k\big)\,d\mathbf x
\;=\;A_k\Big[\Phi\big(\tfrac{i+\frac12-\mu_{y}}{\sigma_{y}}\big)-\Phi\big(\tfrac{i-\frac12-\mu_{y}}{\sigma_{y}}\big)\Big]
\Big[\Phi\big(\tfrac{j+\frac12-\mu_{x}}{\sigma_{x}}\big)-\Phi\big(\tfrac{j-\frac12-\mu_{x}}{\sigma_{x}}\big)\Big],
\end{equation}
the product of two one-dimensional pixel integrals (error functions). The product form is exact only for an
uncorrelated cloud; it is used for clouds with $|\rho|=|\sigma_{xy}|/\sqrt{\sigma_{yy}\sigma_{xx}}<0.5$ (all X-rays),
and the strongly correlated clouds (pieces of tilted tracks, which are broad along the track) keep the Gaussian with
$+\tfrac1{12}$. The footprint is unchanged ($4\sigma$ ellipse, at least 2\,px).

\paragraph{A consistent M-step.} The weighted moments of a cloud's share are now the mean and variance of the
\emph{pixelised} Gaussian, which depend on $\sigma$ and on the sub-pixel position of the centre (a $\sigma=0.2$\,px cloud
centred on a pixel has variance $0.01$\,px$^2$, centred on a pixel edge $0.25$\,px$^2$). The M-step inverts them
with tables computed once (variance against $\sigma^2$ and the centre offset; mean against offset and $\sigma^2$):
$\sigma^2_{y}=V^{-1}\big(\mathrm{var}_y\,|\,\mathrm{frac}(\mu_y)\big)$, $\mu_y=\bar M^{-1}\big(\bar y\,|\,\sigma_y^2\big)$, the
same for $x$. Without this step the exact model and the moment update disagree, and sharp clouds drift in width.

\paragraph{Cost.} The compiled EM pass (section~\ref{sec:phase1}) has an exact kernel (\code{math.erf}); it agrees with
the numpy version to $10^{-15}$\,keV. The time per image increases by $\approx20\,\%$ (the prototype without the compiled
kernel took twice as long).

\paragraph{Two settings that change with it.} (i) The island arbitration (section~\ref{sec:heart}, tracks against several
X-rays) compares the BIC of fixed-shape clouds with free clouds; the exact model fits sharp X-rays with free clouds
better, so the margin was recalibrated on the simulation below from 30 to 100. (ii) With the better fit a sharp
X-ray is one cloud instead of three, and a faint fragment (a few electrons in the tail of a large deposit, fitted as a thin
``track'' cloud) could then join both the X-ray and the deposit and merge them. A fragment --- a cloud below 50\,\el{} and below 10\,\% of its
brightest merge partner --- now joins only that partner's cluster: it can be part of a cluster but cannot connect two.
(The first version used only the 10\,\% condition; next to a 600-keV deposit a 5000-\el{} piece of the same deposit then
counted as faint and the deposit was split, $+25\,\%$ clusters on the crowded CCD.)

\paragraph{The reported width.} The width that separates front- and back-side events is the second moment of the
cluster's shared charge minus $\tfrac1{12}$ (audit D2). It is unbiased for $\sigma\ge0.4$\,px, but below it the pixelised
variance is so small that the noise and dark charge in the footprint dominate it; inverting it exactly does not help
(0.28 for a 0.2-px X-ray). The exact-pixel EM width is unbiased there, so a single cloud with an EM width below 0.4\,px reports
it. Single 5.9-keV X-rays (80 per width, 0.15\,\el{} noise):

\begin{table}[H]\centering\footnotesize
\begin{tabular}{lcccccc}
\toprule
true $\sigma$ [px] & 0.2 & 0.3 & 0.4 & 0.6 & 1.0 & 1.5 \\
\midrule
moment $-\tfrac1{12}$ (before) & 0.168 & 0.289 & 0.400 & 0.600 & 0.997 & 1.505 \\
reported now & 0.250 & 0.304 & 0.401 & 0.600 & 0.997 & 1.505 \\
$E_\mathrm{rec}/E_\mathrm{true}$ (both) & 1.000 & 1.000 & 1.000 & 1.000 & 1.000 & 1.000 \\
\bottomrule
\end{tabular}
\end{table}
\noindent A 0.2-px cloud puts 99\,\% of its charge in one pixel when centred, so its width is barely measurable. Either
way it is far below the front/back cut at 1\,px.

\paragraph{Simulation: sharp X-rays next to diffuse ones.} Skipper blanks (0.1\,\el/pix, 0.16\,\el{} noise), 5.9-keV
X-rays in four scenes, 777 of each (three images):

\begin{table}[H]\centering\footnotesize
\begin{tabular}{llcccc}
\toprule
scene & cluster & \multicolumn{2}{c}{found} & \multicolumn{2}{c}{$|E/E_\mathrm{true}-1|<5\,\%$} \\
 & & before & exact & before & exact \\
\midrule
isolated sharp ($\sigma$ 0.15--0.35) & sharp & 1.000 & 1.000 & 1.000 & 1.000 \\
sharp 3\,px from diffuse ($\sigma$ 1.2--1.6) & sharp & 0.781 & 0.988 & 0.350 & 0.967 \\
 & diffuse & 0.773 & 0.981 & 0.349 & 0.967 \\
two sharp 3\,px apart & sharp & 0.942 & 0.934 & 0.829 & 0.851 \\
isolated diffuse ($\sigma$ 1.0--1.6) & diffuse & 1.000 & 1.000 & 1.000 & 1.000 \\
\bottomrule
\end{tabular}
\end{table}
\noindent The energies of the separated pairs are within $\pm0.8\,\%$ (median). Isolated clusters do not change.

\paragraph{Real data.} The wide-test subset (every 12th Fe-55 image-CCD, every 6th Am-241; 80 paired image-CCDs, the
earlier injection placement), current model against exact pixel integrals:

\begin{table}[H]\centering\footnotesize
\begin{tabular}{lcc}
\toprule
 & before & exact \\
\midrule
Fe-55: injected X-ray pairs 4\,px apart resolved & 0.907 & 0.941 \\
Fe-55: injected X-rays / muons / worms correct & 0.956 / 0.640 / 0.839 & 0.958 / 0.640 / 0.839 \\
Am-241: injected pairs resolved & 0.607 & 0.655 \\
real front-side K$\alpha$ (exact $-$ before, median over image-CCDs) & \multicolumn{2}{c}{$-0.2$\,eV (68\,\%: $-1.8$ to $+2.8$)} \\
K$\alpha$ clusters: $\chi^2$ flag median (90th percentile) & 6.46 (33.3) & 1.31 (4.18) \\
K$\alpha$ clusters: front-side ($\sigma_{xy}<1$) fraction & 0.591 & 0.587 \\
time per image-CCD, Fe-55 / Am-241 (median) & 0.8 / 42\,s & 0.9 / 62\,s \\
\bottomrule
\end{tabular}
\end{table}
\noindent The model now describes the X-rays it fits (the $\chi^2$ flag falls five-fold), so that flag separates
contaminated clusters better; its cut has to be recalibrated. The real front-side K$\alpha$ does not move.

\paragraph{The crowded CCD.} The 29 Am-241 images of the crowded CCD (AMtest EXT4) with the constrained line fit of
section~\ref{sec:perf} (\AGMX, $\sigma_{xy}<0.7$, blend $<0.02$):

\begin{table}[H]\centering\footnotesize
\begin{tabular}{lcccc}
\toprule
 & clusters (all) & $k$ & dev/ndf & $Z$ (Np L, $\gamma$\,26.3) \\
\midrule
published (section~\ref{sec:perf}) & --- & 1.0014 & 1.98 & 34.3 \\
before (with the hot-column fix below) & 497\,765 & $1.0016\pm0.0005$ & 1.75 & 32.8 \\
exact pixel integrals & 504\,440 & $1.0006\pm0.0003$ & 2.59 & 37.4 \\
\quad + ped $<10$ & & $0.9990\pm0.0003$ & 2.12 & 29.9 (before 26.8) \\
\bottomrule
\end{tabular}
\end{table}
\noindent The lines are more significant (L$\beta$ $19.9\to22.4\,\sigma$, L$\gamma$ $17.9\to20.6$, $\gamma$\,26.3
$18.3\to21.1$) and the energy scale is closer to 1, but the deviance per degree of freedom of the spectrum is higher
(1.75 $\to$ 2.59): the shape between the lines is described less well. This is not understood yet and is an open item.

\paragraph{A hot-column bug found on the way.} On this CCD the hot-column correction of the audit (B1) flagged 2400 of
the 3348 columns as hot, with $\approx17$\,\el/pix: its clip was iterated with a threshold that grew with the excess it
measured, so on a charge sheet it swallowed the event charge. Every energy was pulled down by $\approx0.5\,\%$ (59.5-keV
peak 59.55 $\to$ 59.25\,keV on one image; $k=0.975$ on the 29 images). A hot column is now a single column standing
out from its $\pm10$ neighbours (one clip around the column median), and nothing is flagged when more than 10\,\% of the
columns would be; on this CCD no column is flagged, and the published numbers are recovered (table above, ``before'').

\paragraph{EM convergence.} The EM runs a fixed 30 iterations (section~\ref{sec:heart}). Running it to convergence
(largest centre step $<0.01$\,px, at most 300 iterations; 99.4\,\% of the fits converged) was tested with \AGMX{} on the
29 crowded images (before the hot-column fix, so $k$ is low in both): $Z=6.2\to4.3$ with the blend cut and $6.0\to7.2$
with the ped cut in addition, dev/ndf $1.28\to1.04$ and $1.51\to1.41$. There is no consistent gain, and the cluster
energies were already stable at 30 iterations (Step~2). \textbf{The fixed 30 iterations are kept.}


\paragraph{Switch.} \code{AGMX\_PIX\_EXACT=0} restores the earlier pixel model, island margin 30 and merge graph.
